\section*{अध्याय \ref{ch:b1:water-small-molecules} --- जल और छोटे जैव-अणु}

\begin{solution}{exo:b1:water-small-molecules:1}
इस मुड़े हुए, ध्रुवी अणु की ऑक्सीजन अंशतः ऋणात्मक और हाइड्रोजन अंशतः
धनात्मक हैं: वह किसी धनायन को अपनी ऑक्सीजनों से और किसी ऋणायन को अपनी
हाइड्रोजनों से घेर लेता है, उन्हें स्थायी करता हुआ (जलयोजन), और वह ध्रुवी
समूहों से \omterm{def:b1:water-small-molecules:water}{हाइड्रोजन बंध} बनाता है। तेल में बँधने लायक कोई आवेश या ध्रुवी
समूह नहीं; उसके चारों ओर जल-अणुओं को एक पिंजरे में क्रमबद्ध होना पड़ता है,
जो प्रतिकूल है, इसलिए तेल बाहर कर दिया जाता है।
\end{solution}

\begin{solution}{exo:b1:water-small-molecules:2}
\omterm{def:b1:water-small-molecules:water}{हाइड्रोजन बंध} (\qtyrange{4}{20}{kJ/mol}): DNA में क्षारक-युग्मन, प्रोटीन
की द्वितीयक संरचना। आयनी बंध (\omterm{def:b1:water-small-molecules:water}{जल} में \qty{12}{kJ/mol}): प्रोटीनों में
लवण-सेतु, क्रियाधार का बँधना। जलविरोधी प्रभाव (कोई बंध नहीं,
पर कुल मिलाकर तुलनीय): झिल्ली का जुड़ना, प्रोटीन वलन। \omterm{def:b1:water-small-molecules:bonds}{वान डर
वाल्स} (हर एक \qtyrange{1}{4}{kJ/mol}): किन्हीं भी दो
पूरक पृष्ठों का कसकर जमना, एंजाइम--क्रियाधार का बैठना।
\end{solution}

\begin{solution}{exo:b1:water-small-molecules:3}
$\mathrm{pH} = -\log(4\times 10^{-8}) = 7.4$। \omterm{def:b1:water-small-molecules:ph}{pH} 1.5 पर,
$[\mathrm{H^+}] = 10^{-1.5} = \qty{0.032}{mol/L}$।
\end{solution}

\begin{solution}{exo:b1:water-small-molecules:4}
लगभग 5.8 से 7.8 तक (p$K_a \pm 1$)। 7.4 पर: $\log(\text{क्षार}/\text{अम्ल})
= 0.6$, अनुपात $4$।
\end{solution}

\begin{solution}{exo:b1:water-small-molecules:5}
\omterm{def:b1:water-small-molecules:ph}{pH} 4.76 पर: \qty{50}{\%}। 7.2 पर: $\log(\mathrm{A^-}/\mathrm{HA}) =
2.44$, अनुपात $275$: \qty{99.6}{\%} आयनित। \omterm{def:b1:water-small-molecules:ph}{pH} 2 पर: अनुपात $10^{-2.76} =
1/575$: \qty{0.2}{\%} आयनित। उदासीन अम्ल-रूप \omterm{def:b1:membranes-transport:membrane}{लिपिड
द्विस्तर} पार करता है; ऐसीटिक अम्ल आमाशय में सोखा जाता है, जहाँ वह
अनआयनित रहता है, और \omterm{def:b1:cell-unit-of-life:cell}{कोशिकाद्रव्य} में ऐसीटेट के रूप में फँस जाता है।
\end{solution}

\begin{solution}{exo:b1:water-small-molecules:6}
\omterm{def:b1:water-small-molecules:ph}{pH} $= 6.8 + \log 1 = 6.8$। HCl के बाद: क्षार $40$, अम्ल $60$: \omterm{def:b1:water-small-molecules:ph}{pH} $= 6.8 +
\log(40/60) = 6.62$। शुद्ध \omterm{def:b1:water-small-molecules:water}{जल} में, $[\mathrm{H^+}] = \qty{0.01}{mol/L}$:
\omterm{def:b1:water-small-molecules:ph}{pH} $2$।
\end{solution}

\begin{solution}{exo:b1:water-small-molecules:7}
$500/2.4 = \qty{208}{g}$। किसी अणु को वाष्पित करने के लिए उसके चारों
\omterm{def:b1:water-small-molecules:water}{हाइड्रोजन बंध} तोड़ने पड़ते हैं; \qty{2.4}{kJ/g} यानी \qty{43}{kJ/mol},
मोटे तौर पर प्रति अणु दो बंधों की ऊर्जा (हर बंध दो में बँटा होता है)।
\end{solution}

\begin{solution}{exo:b1:water-small-molecules:8}
$\Delta G^{\circ\prime} = -RT\ln 19 = -2.48\times 2.94 =
\qty{-7.3}{kJ/mol}$। \omterm{def:b1:cell-unit-of-life:cell}{कोशिका} में: $\Delta G = -7.3 + 2.48\ln(2/0.02) =
-7.3 + 11.4 = \qty{+4.1}{kJ/mol}$: अभिक्रिया उलटी दिशा में,
ग्लूकोज-1-फ़ॉस्फ़ेट की ओर चलती है (जैसा वह ग्लाइकोजन-संश्लेषण में करती है)।
\end{solution}

\begin{solution}{exo:b1:water-small-molecules:9}
\omterm{def:b1:water-small-molecules:water}{जल} \qty{4}{\celsius} पर सबसे सघन होता है; उससे ठंडा \omterm{def:b1:water-small-molecules:water}{जल} और बर्फ़ हलके होते
हैं और पृष्ठ पर रह जाते हैं, इसलिए बर्फ़ ऊपर बनती है और नीचे के तरल को
कुचालित कर देती है, जो सारा जाड़ा \qty{4}{\celsius} के पास बना रहता है।
यदि बर्फ़ डूबती, तो झीलें तली से ऊपर तक ठोस जम जातीं और गरमी में केवल
पृष्ठ पर पिघलतीं; किसी जाड़े भर मीठे \omterm{def:b1:water-small-molecules:water}{जल} का जीवन असंभव
हो जाता।
\end{solution}

\begin{solution}{exo:b1:water-small-molecules:10}
$RT = \qty{2.58}{kJ/mol}$। विश्राम में: $Q = 0.9\times 10^{-3}\times 8\times
10^{-3}/8\times 10^{-3} = 9\times 10^{-4}$, $\ln Q = -7.0$: $\Delta G =
-30.5 - 18.1 = \qty{-48.6}{kJ/mol}$। थकी हुई: $Q = 3\times 10^{-3}\times
30\times 10^{-3}/5\times 10^{-3} = 0.018$, $\ln Q = -4.0$: $\Delta G =
-30.5 - 10.4 = \qty{-40.9}{kJ/mol}$। थकान ने प्रति \omterm{def:b1:water-small-molecules:atp}{ATP} ऊर्जा
\qty{16}{\%} काट दी है, जो उन पंपों और प्रेरकों को धीमा करने के लिए बहुत
है जिन्हें उसका अधिकांश चाहिए।
\end{solution}

\begin{solution}{exo:b1:water-small-molecules:11}
$[\mathrm{H^+}] = 10^{-7.2} = \qty{6.3e-8}{mol/L}$; आयतन
\qty{1e-12}{L}: $\qty{6.3e-20}{mol}\times\num{6e23} = 38$ प्रोटॉन। और
$\num{3e9}$ प्रतिरोधी समूहों के सामने: ``\omterm{def:b1:water-small-molecules:ph}{pH}'' प्रोटॉनों की गिनती नहीं बल्कि
प्रतिरोधकों की प्रोटॉनन-अवस्था है, जो \omterm{def:b1:water-small-molecules:water}{जल} से प्रोटॉन उन थोड़े मुक्त
प्रोटॉनों के मायने रखने से एक अरब गुना तेज़ बदलते रहते हैं।
\end{solution}

\begin{solution}{exo:b1:water-small-molecules:12}
साम्य पर \omterm{def:b1:water-small-molecules:atp}{ATP}, ADP और फ़ॉस्फ़ेट लगभग
$10^{-5}$ के अनुपात पर बैठते; \omterm{def:b1:cell-unit-of-life:cell}{कोशिका} \omterm{def:b1:water-small-molecules:atp}{ATP} को ADP से हज़ार गुना ऊपर रखती है,
ताकि उसका \omterm{prop:b1:water-small-molecules:families}{जल-अपघटन} कुछ नहीं के बजाय \qty{50}{kJ/mol} छोड़े। हर चढ़ाई वाली
प्रक्रिया --- संश्लेषण, परिवहन, गति --- उसी विस्थापन से युग्मित है;
अपचय की एकमात्र भूमिका उसे बनाए रखना है। जिस \omterm{def:b1:cell-unit-of-life:cell}{कोशिका} का \omterm{def:b1:water-small-molecules:atp}{ATP}
साम्य पर आ जाए उसके पास खर्च करने को कोई प्रवणता नहीं बचती: पंप रुक जाते
हैं, आयन रिसते हैं, \omterm{def:b1:cell-unit-of-life:cell}{कोशिका} मिनटों में फूलकर मर जाती है। जीवित होना उस
असाम्य को बनाए रखना है, कोई पदार्थ नहीं।
\end{solution}

\begin{solution}{pb:b1:water-small-molecules:1}
\textbf{1.} $6.1 + \log(24/1.2) = 6.1 + 1.301 = 7.40$।
\textbf{2.} $10^{-7.4} = \qty{4.0e-8}{mol/L} = \qty{40}{nmol/L}$।
\textbf{3.} $40\times 5 = \qty{200}{nmol}$, यानी बाइकार्बोनेट से छह लाख
गुना कम।
\textbf{4.} $\mathrm{CO_2} = 1.2\times 46/40 = \qty{1.38}{mmol/L}$;
\omterm{def:b1:water-small-molecules:ph}{pH} $= 6.1 + \log(25/1.38) = 6.1 + 1.258 = 7.36$।
\textbf{5.} किसी फ्लास्क में $20:1$ का अनुपात क्षार सोखने के लिए बहुत कम
अम्ल-साथी छोड़ता है, और p$K_a$ से एक इकाई दूर क्षमता नीची रहती है;
शरीर में अम्ल-साथी एक गैस है जिसकी सांद्रता फेफड़े तय करते हैं और बिना
सीमा फिर से बना सकते हैं, इसलिए वह अनुपात p$K_a$ की परवाह किए बिना
संवातन से फिर बैठाया जा सकता है।
\textbf{6.} $24/1.2 = 20$; क्षार जोड़ी का $20/21 = \qty{95}{\%}$ है।
\textbf{7.} $60/5 = \qty{12}{mmol/L}$।
\textbf{8.} बाइकार्बोनेट $24 - 12 = \qty{12}{mmol/L}$; $\mathrm{CO_2}$
$1.2 + 12 = \qty{13.2}{mmol/L}$।
\textbf{9.} \omterm{def:b1:water-small-molecules:ph}{pH} $= 6.1 + \log(12/13.2) = 6.1 - 0.04 = 6.06$।
\textbf{10.} $[\mathrm{H^+}] = \qty{0.012}{mol/L}$: \omterm{def:b1:water-small-molecules:ph}{pH} $1.92$।
\textbf{11.} $1.34$ इकाइयों की खिसकन: वह रक्त को \omterm{def:b1:water-small-molecules:ph}{pH} 1.9 से बचा लेता है
पर 6.06 उससे बहुत नीचे है जितना कोई एंजाइम सह सके। किसी संवृत
निकाय के रूप में, खराब।
\textbf{12.} $[\mathrm{H^+}]$ \qty{40}{nmol/L} से चढ़कर $10^{-6.06} = \qty{870}{nmol/L}$ हो जाता है: यानी जोड़े गए \qty{12}{mmol/L} के लिए \qty{0.83}{\micro mol/L} की बढ़त, यानी \num{14000} में एक प्रोटॉन मुक्त रह जाता है; बाकी बाइकार्बोनेट पर हैं।
\textbf{13.} \omterm{def:b1:water-small-molecules:ph}{pH} $= 6.1 + \log(12/1.2) = 6.1 + 1.0 = 7.10$।
\textbf{14.} बनी $\mathrm{CO_2}$: $\qty{12}{mmol/L}\times
\qty{5}{L} = \qty{60}{mmol}$।
\textbf{15.} \qty{10}{mmol} के ऊपर एक मिनट में \qty{60}{mmol}:
यानी सात गुना।
\textbf{16.} $\mathrm{CO_2} = \qty{0.9}{mmol/L}$; \omterm{def:b1:water-small-molecules:ph}{pH} $= 6.1 +
\log(12/0.9) = 6.1 + 1.125 = 7.22$।
\textbf{17.} 7.4 के पास विवृत बाइकार्बोनेट: \omterm{def:b1:water-small-molecules:ph}{pH} एक इकाई तब बदलता है
जब बाइकार्बोनेट दस गुना गिरे, \qty{24}{} से
\qty{2.4}{mmol/L} तक, यानी प्रति इकाई \qty{21.6}{mmol/L} अम्ल; पर
पहली इकाइयों भर स्थानीय ढाल $2.3\times[\mathrm{HCO_3^-}] =
\qty{55}{mmol/L}$ प्रति इकाई है। प्रोटीनों के \qty{25}{mmol/L} प्रति
इकाई के सामने, बाइकार्बोनेट प्रोटॉनों का लगभग $55/80 = \qty{69}{\%}$ लेता
है, और प्रोटीन \qty{31}{\%}।
\textbf{18.} बाइकार्बोनेट $0.69\times 12 = \qty{8.3}{mmol/L}$ लेता है:
$24 - 8.3 = \qty{15.7}{mmol/L}$; \omterm{def:b1:water-small-molecules:ph}{pH} $= 6.1 + \log(15.7/1.2) = 6.1 +
1.117 = 7.22$ (प्रोटीनों का हिस्सा \qty{25}{mmol/L}
प्रति इकाई $\times 0.18$ इकाई $= \qty{4.5}{mmol/L}$ के अनुकूल, जो उनके
\qty{3.7}{mmol/L} के पास है)।
\textbf{19.} अम्ल हटाने से वह बाइकार्बोनेट फिर बन जाता है जो खर्च हुआ था
(\qty{12}{mmol/L} से वापस \qty{24}{}) और, $\mathrm{CO_2}$ थामे रखने पर,
\omterm{def:b1:water-small-molecules:ph}{pH} 7.40 पर लौट आता है।
\textbf{20.} प्रति दिन $60/5 = \qty{12}{mmol/L}$ बाइकार्बोनेट खोता है;
तीन दिन बाद $24 - 36 < 0$: बाइकार्बोनेट चुक जाता है और \omterm{def:b1:water-small-molecules:ph}{pH}
7 से कहीं नीचे ढह जाता है --- व्यवहार में रोगी को बाइकार्बोनेट देना या
अपोहन करना ही पड़ता है; और अकेले एक दिन बाद, $6.1 + \log(12/1.2) =
7.10$।
\textbf{21.} $\mathrm{CO_2} = 1.2\times 60/40 = \qty{1.8}{mmol/L}$;
\omterm{def:b1:water-small-molecules:ph}{pH} $= 6.1 + \log(24/1.8) = 6.1 + 1.125 = 7.22$। \qty{33}{} पर
बाइकार्बोनेट के साथ: $6.1 + \log(33/1.8) = 6.1 + 1.263 = 7.36$।
\textbf{22.} फेफड़े कुछ ही साँसों में हर बदल देते हैं (घुली
$\mathrm{CO_2}$), क्योंकि वह एक गैस है जिसे वे बाहर छोड़ते हैं; और
वृक्क अंश बदलते हैं (बाइकार्बोनेट), अम्ल उत्सर्जित करके और बाइकार्बोनेट
फिर से बनाकर, जो घंटों से दिनों की प्रक्रिया है।
\textbf{23.} \omterm{def:b1:water-small-molecules:ph}{pH} 5 पर, $[\mathrm{H^+}] = \qty{10}{\micro mol/L}$: यानी \qty{1.5}{L} में \qty{15}{\micro mol} मुक्त, जो \qty{60}{mmol} का चार-हज़ारवाँ भाग है; अम्ल प्रतिरोधकों से बँधकर निकलता है --- अमोनियम के रूप में और डाइहाइड्रोजन फ़ॉस्फ़ेट के रूप में।
\textbf{24.} अम्ल खोने से \omterm{def:b1:water-small-molecules:ph}{pH} चढ़ता है: बाइकार्बोनेट
$100/5 = \qty{20}{mmol/L}$ चढ़कर \qty{44}{} हो जाता है; \omterm{def:b1:water-small-molecules:ph}{pH} $= 6.1 + \log(44/1.2) =
6.1 + 1.564 = 7.66$ (क्षारमयता), इससे पहले कि फेफड़े और वृक्क
भरपाई करें।
\textbf{25.} संवृत: 7.40 से 6.06 तक, यानी $-1.34$ की खिसकन। विवृत,
अकेले बाइकार्बोनेट से: 7.10 तक, यानी $-0.30$ की खिसकन। प्रोटीनों के साथ
विवृत: 7.22 तक, यानी $-0.18$ की खिसकन।
\end{solution}
